%!TEX root = ../main.tex
\begin{center}
    \large

    \begin{singlespace}
        \textbf{\englishTitle{}} \\[0.5cm]
    \end{singlespace}

    \begin{singlespace}
        Student : \studentEnName{}  \hspace{1.0cm}
        % 兩個指導教授要寫 Advisors
        \ifdefined\advisorCnNameB
            Advisors: Dr.\, \advisorEnName \\
            \hspace{6.6cm} Dr.\, \advisorEnNameB  \\
        \else
            Advisor: Dr.\, \advisorEnName \\
        \fi
    \end{singlespace}

    \vspace{0.5cm}
    \begin{singlespace}
        \NameofDepartmentInstituteEN\\
        National Yang Ming Chiao Tung University\\[0.2cm]
    \end{singlespace}

    \textbf{Abstract} \\[0.5cm]

\end{center}

\normalsize

Deep-learning-based automatic modulation classifiers (AMC) achieve high
accuracy under benign channel conditions but are vulnerable to adversarial
perturbations. Optimization-based attacks such as the Carlini--Wagner ($L_2$)
attack (CW) and the Elastic-Net Attack on Deep Neural Networks (EAD) can
substantially reduce the accuracy of an Adaptive Wavelet Network (AWN)
classifier on RML2016.10a, even when the underlying waveform remains largely
decodable. This thesis proposes \emph{Adaptive-K}, a receiver-side
frequency-domain input-transformation defense that recovers classification
accuracy without retraining the classifier.

Adaptive-K combines three mechanisms in a single shared-FFT pass:
spectral-flatness routing for wideband signals, an SNR-adaptive K-cap that
balances perturbation suppression against signal preservation, and per-sample
magnitude-knee estimation to select the number of retained spectral
components. The method preserves the dominant modulation-carrying spectral
components while suppressing non-dominant attack-induced spectral artifacts.

We evaluate Adaptive-K against five attacks---CW, EAD-L1, EAD-EN, FGSM, and
PGD---on RML2016.10a using a pretrained AWN classifier. We compare against
five per-SNR-calibrated classical filters (Kalman, Wiener, Savitzky-Golay,
Gaussian, FIR) and an empirical randomized-smoothing baseline
($k = 20$ noisy copies, $\sigma = 0.01$; no formal certificate). Across SNR $\geq 0$\,dB, CW reduces weighted-average accuracy to
35.4\% and EAD-L1 to 6.6\%. Adaptive-K recovers CW accuracy to 71.3\% and
EAD-L1 accuracy to 70.8\%, outperforming the strongest classical baseline by
6.2 and 8.5 percentage points respectively. On PGD it improves accuracy from
32.6\% to 43.4\%. Clean accuracy degrades by less than 0.5 percentage points,
and the defense adds less than 0.1\,ms latency per batch of 64 samples in our
GPU implementation.

These results show that adaptive spectral filtering is an effective
lightweight defense for recovering AMC performance under defense-unaware
adversarial attacks. The current evaluation does not claim robustness against
adaptive defense-aware attacks; evaluating such attacks remains an important
direction for future work.

\vspace{1cm}

% 5-7 Keywords (English)
\textbf{Keywords: automatic modulation classification, adversarial robustness,
spectral defense, FFT Top-$K$, RML2016.10a, signal processing.}
